\section{Kết luận}
\label{sec:conclusion}

Chuyên đề đã nghiên cứu và hiện thực hóa một \emph{cơ chế phát hiện tấn công leo
thang đặc quyền} cho môi trường multi-tenant Kubernetes dựa trên phân tích lời gọi
hệ thống. Cơ chế kế thừa thiết kế lõi phát hiện bất thường SyscallAD (VAE + Rừng cô
lập) của công trình tham chiếu DeSFAM~\cite{desfam_original}, định vị các đặc trưng
syscall đặc quyền làm tín hiệu cho leo thang đặc quyền; huấn luyện trên bộ dữ liệu
công khai DongTing và được triển khai thành một bộ phát hiện chạy trực tiếp bằng
Tetragon/eBPF giám sát theo namespace trên Kubernetes, toàn bộ đóng gói bằng Docker
để bảo đảm tái lập.

Các kết quả chính như sau. \emph{Về phát hiện ngoại tuyến}, lõi phát hiện đạt định
hướng và khả năng xếp hạng tốt (VAE AUC kiểm định $0{,}944$, AP tập hợp $0{,}990$,
precision $0{,}999$) nhưng còn khác biệt về F1/recall so với mức công bố do ngưỡng
ưu tiên precision. \emph{Về triển khai trực tiếp}, mô hình giữ được khả năng phân
tách khi phục vụ trực tiếp (lành tính $\le 0{,}19$ so với các kịch bản leo thang
đặc quyền $0{,}278$--$0{,}930$), cho phép ngưỡng vận hành $T_0=0{,}25$. \emph{Về
điều kiện thành công}, yếu tố then chốt là căn chỉnh không gian đặc trưng (bảng chữ
cái cảm biến 23 syscall) giữa huấn luyện và phục vụ, được ràng buộc bằng kiểm thử
ngang hàng.

Tóm lại, báo cáo trình bày một cơ chế phát hiện khả thi cho leo thang đặc quyền
trong multi-tenant Kubernetes, kèm một đối chiếu \emph{trung thực} với công trình
tham chiếu: chỉ rõ điều gì tương đồng, điều gì khác biệt và vì sao, đồng thời
lấp khoảng cách ``mô phỏng $\to$ triển khai'' bằng một bộ phát hiện trực tiếp. Hướng
tiếp theo gồm chọn ngưỡng tối ưu F1, hiện thực hóa khâu chặn trong nhân (Giai đoạn~3)
để \emph{ngăn chặn} chứ không chỉ phát hiện, mở rộng khối lượng công việc đa tenant
và dùng khai thác CVE thực trong môi trường cách ly.
